\section{结论}
\label{sec:conclusion}

我们提出 \ProjectName，一个把 harness 视为模型与环境之间一等接口的可组合运行时铸造厂。该接口可由带类型原语组合，可依据执行轨迹演化，也可在统一改进循环中与模型训练相结合。在五个基准和三个模型家族上，\ProjectName 通过组合底座上的轨迹驱动演化，取得最高 $+44.0\%$ 的增益（15 个配置平均 $+14.5\%$）；在两个基准上，协同演化又比仅演化 harness 进一步提升 $+4.7\%$。这些结果表明，智能体进步未必只依赖模型扩展：根据执行反馈组合和演化运行时接口，是一条互补且可操作的路径；对于能力受限、harness 层增益最大的智能体而言尤其如此。
